% ============================================================
% Section 1: Introduction
% ============================================================
\section{Introduction}
\label{sec:introduction}

Physically Unclonable Functions (PUFs) are hardware security primitives that
derive device-unique secrets from intrinsic manufacturing variations.  Among
strong PUFs, the XOR Arbiter PUF (XOR APUF) has attracted the most attention
from the modeling attack community: an attacker collects challenge--response
pairs (CRPs) from a device and trains a machine learning model to predict
future responses.  The conventional wisdom holds that increasing the number
of parallel APUFs $k$ renders the PUF exponentially harder to model, and
that a sufficiently complex model---deep neural network, support vector
machine, evolution strategy---will eventually succeed given enough data.

This binary framing---``can $k$-XOR be broken?''---has guided a decade of
research, but it overlooks the most interesting aspect: ML-based modeling
attacks on XOR APUFs exhibit a \emph{probabilistic phase transition} in which,
for each $k$, the probability of attack success $P(\text{success} \mid N, k)$
follows an S-shaped curve as a function of CRP budget $N$.  Below the
transition region, \emph{every} training run yields a failed classifier:
for $k \geq 6$, failure means accuracy indistinguishable from random guessing
($\approx 50\%$); at $k = 4$, failures show intermediate accuracies
(58--76\%).  Above the transition, \emph{every} run succeeds ($> 99\%$).
In the transition region, identical conditions with different random seeds
can produce either outcome.  The transition is not a sharp threshold but a
gradual S-curve whose absolute width (in CRPs) increases with $k$.

\paragraph{Contributions.}
This paper makes the following contributions:

\begin{itemize}
    \item \textbf{Probabilistic S-curve characterization.}  We systematically
    measure $P(\text{success} \mid N, k)$ for $k = 2, 4, 6, 8$ at
    multiple data levels, with 3--30 independent training runs per
    configuration (\S\ref{sec:success_curves}).  The transition midpoint
    $N_{50}$ scales rapidly: from $\sim 10^5$ CRPs for 2-XOR to
    $\sim 8.0 \times 10^6$ CRPs for 8-XOR (logistic fit), and
    $> 1.5 \times 10^7$ CRPs for 9-XOR (preliminary lower bound).

    \item \textbf{Architecture independence across model families and
    optimizers.}  Eight distinct MLP
    architectures---varying in depth, width, activation, residual
    connections, and parameter count by $15\times$---and two
    evolution-strategy baselines (PSO, CMA-ES) with the structurally correct
    delay parameterization all fail identically
    below the data threshold, yet the same model succeeds given sufficient
    data: data quantity, not model design or optimization paradigm, is the
    bottleneck (\S\ref{sec:arch_independence}, \S\ref{sec:evo_baselines}).

    \item \textbf{Linear model impossibility.}  Logistic regression produces
    random accuracy ($\approx 50\%$) for all $k \geq 4$ regardless of data
    volume, confirming that the XOR nonlinearity is fundamentally
    inaccessible to linear classifiers (\S\ref{sec:lr_bottleneck}).

    \item \textbf{iPUF vulnerability.}  Every Interpose PUF configuration
    tested, up to $(4,4)$-iPUF, succumbs to the split attack with
    1{,}000{,}000 CRPs and $> 99\%$ accuracy.  Going beyond the established
    breakability result~\cite{wisiol2020splitting}, we quantify the
    effective security level relative to our XOR APUF scaling law: the split
    attack reduces the equivalent XOR complexity by exactly 2, yielding a
    CRP-budget gap that grows from $\approx$5$\times$ at $(2,4)$ to
    $\approx$21$\times$ at $(4,4)$ (\S\ref{sec:ipuf}).

    \item \textbf{Practical bound for 8-XOR.}  An 8-XOR APUF with CRP
    exposure limited to $5 \times 10^6$ CRPs provides practical security
    against single-GPU ML attacks: the attacker succeeds with at most 28\%
    probability.  For 9-XOR, the $N_{50}$ estimate of $>1.5\times10^7$ CRPs
    already exceeds the memory of our 12~GB GPU; at $k = 10$, no successes were observed at
    $10^7$ CRPs, and 15M CRPs cannot be loaded even on a 24~GB GPU.
\end{itemize}

\paragraph{Related work.}
Modeling attacks on XOR APUFs date back to the earliest PUF
studies~\cite{ruhrmair2010modeling}.  Recent work has explored evolution
strategies~\cite{becker2015cia, calypso2024}, deep
learning~\cite{fei2024mope, wisiol2020splitting}, and chosen-challenge
attacks~\cite{sayadi2025}.  The split attack of Wisiol~et~al.~\cite{wisiol2020splitting}
is the current state of the art for Interpose PUFs.  To our knowledge,
however, no prior work has systematically characterized the \emph{probability}
of attack success as a function of data volume across multiple XOR
complexities, nor identified the probabilistic phase transition as a
governing principle.

\paragraph{Organization.}
Section~\ref{sec:background} provides background on XOR APUFs and ML-based
attacks.  Section~\ref{sec:setup} describes our experimental methodology.
Section~\ref{sec:results} presents the experimental results across four
dimensions.  Section~\ref{sec:discussion} interprets the findings, discusses
their implications for PUF security evaluation, and notes limitations.
Section~\ref{sec:conclusion} concludes.
